\documentclass[11pt]{article}
\usepackage[margin=1in]{geometry}
\usepackage{amsmath,amssymb,graphicx,booktabs,hyperref,xcolor}
\usepackage{siunitx}
\hypersetup{colorlinks=true,linkcolor=blue,urlcolor=blue}
\title{Replication Report: \\ \emph{FeGe as a building block for the kagome
1:1, 1:6:6, and 1:3:5 families} \\ \normalsize (Jiang, Hu, C\u{a}lug\u{a}ru,
Felser, Blanco-Canosa, Weng, Xu, Bernevig --- arXiv:2311.09290v2)}
\author{TEXTURES-100 automated replication (loop-current class)}
\date{2026-07-19}
\begin{document}
\maketitle

\begin{abstract}
We replicate the machine-checkable core of arXiv:2311.09290, a d-orbital
tight-binding + $S$-matrix flat-band framework for FeGe-class kagome materials.
\textbf{Scope caveat:} this paper was filed under the \emph{loop-current} class,
but it contains \emph{no} loop-current / TRS-breaking-flux physics. The only
overlap with the supplied reusable kagome loop-current kernel is the shared
nearest-neighbor (NN) kagome tight-binding Bloch Hamiltonian and its flat-band /
Dirac / van-Hove spectrum --- i.e. the flux$=0$ limit of the kernel's
\texttt{KagomeModel}. We reuse that substrate verbatim (with provenance) and add
the paper-specific bipartite-crystalline-lattice (BCL) flat-band-counting
machinery. Five concrete claims were tested with real code; all five reproduce
quantitatively. \textbf{Verdict: reproduced (within scope).} Coverage $6/10$,
Agreement $10/10$.
\end{abstract}

\section{Paper summary}
The authors argue that the ``spaghetti'' Fe-$d$ band manifold of FeGe (space
group 191) decomposes by symmetry and chemistry into three decoupled orbital
groups, $H(k)=H_1\oplus H_2\oplus H_3$, with
$H_1=(d_{xy},d_{x^2-y^2})\oplus\mathrm{Ge}_t(p_x,p_y)$,
$H_2=(d_{xz},d_{yz})\oplus\mathrm{Ge}\,p_z$, $H_3=d_{z^2}$. Each group forms a
\emph{bipartite crystalline lattice} (BCL); the $S$-matrix formalism then
explains the origin of the quasi-flat bands near $E_F$. The bare building block
is the NN kagome Hamiltonian, whose exactly flat band becomes the FeGe
quasi-flat band once realistic (small) further-neighbor and inter-orbital
hoppings are switched on. The paper further gives cRPA interactions with an
approximate hidden $O_h$ symmetry ($U_1{=}\SI{1.41}{eV}$, $U_2{=}\SI{1.22}{eV}$,
onsite $U\approx\SI{4.15}{eV}$) and a Hartree--Fock reproduction of the A-type
AFM order --- these DFT/cRPA-dependent parts are out of scope for an
independent-code replication (see \S\ref{sec:oos}).

\section{Method}
All spectral results reuse the shared kernel
\texttt{loop\_current\_kagome\_kernel.py}
(\texttt{\char`~/Dropbox/XFER/TEXTURES-100/shared-kernels/}), imported directly
so results are provably produced by it. The kernel's \texttt{KagomeModel} with
\texttt{flux=0} is exactly the NN kagome TB $H_0(k)$ used throughout the paper;
in its convention $H_{0,ab}=-2t\cos(\mathbf{k}\cdot\mathbf{a}_i/2)$ the flat band
sits at $E=+2t$ and the two dispersive bands touch at $K$. We add BCL machinery:
for a chiral bipartite Hamiltonian $H(k)=\left(\begin{smallmatrix}0 & S(k)\\
S^\dagger(k) & 0\end{smallmatrix}\right)$ with sublattice orbital counts
$N_L,N_{\tilde L}$, the number of $E{=}0$ flat bands is
$N_L+N_{\tilde L}-2\,\mathrm{rank}(S_k)$ (App.\ D). We reconstruct the paper's
own $S_{p_{txy},d_1}(k)$ block (Eq.\ S2.18) and count numerical zero modes over
the BZ.

\section{Claims tested and results}
\begin{table}[h]\centering
\begin{tabular}{clcc}
\toprule
ID & Claim & Paper & This work \\
\midrule
C1 & NN kagome: flat band energy ($/t$) & $+2$ & $2.000$ (width $7\!\times\!10^{-15}$) \\
C1 & Dirac degeneracy gap at $K$ & $0$ & $1\!\times\!10^{-15}$ \\
C2 & Case \#1 flat bands ($N_d{-}N_p{=}6{-}2$) & $4$ & $4$ \\
C2 & Case \#4 flat bands ($N_{d2}{-}N_p{=}3{-}2$) & $1$ & $1$ \\
C3 & BCL theorem $N_L{+}N_{\tilde L}{-}2\,r$ & $4$ & $4$ (numeric zero modes) \\
C4 & flat-band width, NN only & $0$ & $2.6\!\times\!10^{-15}$~eV \\
C4 & flat-band width, $t^{NNN}_{d1}{=}0.03$ & quasi-flat & $0.205$~eV \\
C4 & flat-band width, $t^{NNN}{=}0.30$ (large) & dispersive & $1.95$~eV \\
C5 & DOS vHS peak energy $=E(M)$ & $E_M$ & $-2.004$ vs $-2.000$ \\
\bottomrule
\end{tabular}
\caption{All nine numerical sub-checks reproduce. Units: $t{=}1$ for C1--C3/C5;
C4 uses the paper's fitted $t^{NN}_{d1}{=}\SI{0.49}{eV}$.}
\end{table}

\begin{figure}[h]\centering
\includegraphics[width=0.48\textwidth]{../work/bands_kagome.png}\hfill
\includegraphics[width=0.48\textwidth]{../work/dos_kagome.png}
\caption{Left: NN kagome band structure along $\Gamma$--$K$--$M$--$\Gamma$ from
the shared kernel; the flat band sits at $E=+2t$ (dashed red) and the two lower
bands touch in a Dirac cone at $K$. Right: DOS showing the $\delta$-like
flat-band peak at $+2t$ and the logarithmic van-Hove peak at the $M$-point
saddle ($E\approx-2t$). Both reproduce the paper's stated spectral structure
(Figs.\ 2, 4, S2.4).}
\end{figure}

\subsection{C1 --- NN kagome spectrum}
The kernel returns a flat band at exactly $+2t$ (bandwidth $7\times10^{-15}$),
degenerate lower bands at $K$ ($E=-t$, gap $10^{-15}$), and an $M$-point saddle
at $E=0$ in the kernel's zero-of-energy convention. This is the textbook kagome
spectrum the paper takes as its building block.

\subsection{C2/C3 --- BCL flat-band counting}
Reconstructing the paper's inter-sublattice $S$-matrix and diagonalizing the
chiral bipartite Hamiltonian over a BZ mesh yields exactly $4$ zero-energy flat
bands for the full $H_1$ chiral limit ($N_d{=}6$ vs $N_p{=}2$; Case \#1) and $1$
for the single-$d$-group / decoupled limit ($N_{d2}{=}3$ vs $N_p{=}2$; Case
\#4). Both match the paper's arithmetic $N_L-N_{\tilde L}$ and the general
theorem $N_L+N_{\tilde L}-2\,\mathrm{rank}(S)$ (here $\mathrm{rank}(S)=2$
generically over the BZ). A randomized-amplitude control of the same rank
structure confirms the count is amplitude-independent, as the theorem requires.

\subsection{C4 --- quasi-flat mechanism}
The paper's key physical claim: the FeGe quasi-flat band is a bare NN-kagome
flat band that acquires a \emph{small} dispersion because the fitted
$H_{d2}/S_{d1,d2}$ hoppings are small. We confirm the NN-only $d$-band is
\emph{exactly} flat (width $2.6\times10^{-15}$), stays narrow (\SI{0.205}{eV})
under the paper's fitted small $t^{NNN}_{d1}{=}\SI{0.03}{eV}$, and becomes fully
dispersive (\SI{1.95}{eV}) once that hopping is made large --- the exact
mechanism argued in \S II\,A and Fig.\ S2.4.

\section{Out-of-scope items}\label{sec:oos}
The following paper results depend on external DFT/cRPA/Wannier inputs that
cannot be regenerated from first principles here without the authors' DFT
pipeline and MLWFs, and are marked out-of-scope (not faked): (i) the specific
fitted hopping tables per orbital group; (ii) the $85\%$/$97\%$ DFT--TB
wavefunction overlaps; (iii) the cRPA interaction matrix (Table I) and its
hidden $O_h$ symmetry; (iv) the Hartree--Fock AFM magnetic moments; (v) the
1:6:6 / 1:3:5 material extensions. We verified the \emph{model-internal},
first-principles-free claims (flat-band counting, kagome spectral structure,
quasi-flat mechanism), which are the transferable theoretical content.

\section{Verdict}
\textbf{Reproduced (within scope).} Every model-internal, code-checkable claim
reproduces to machine precision or to the paper's stated qualitative behavior.
\textbf{Coverage: 6/10} (the DFT/cRPA-dependent half of the paper is
inherently outside an independent-code replication). \textbf{Agreement: 10/10}
on everything tested.
\end{document}
